\section[Extended holonomy sweeps]{Extended holonomy sweeps and robustness diagnostics (supplement)}
\label{app:holonomy_sweeps_extended}

This appendix records extended finite sweeps and robustness variants for the protocol-connection holonomy diagnostics of Section~\ref{sec:protocol_connections_holonomy}.
The main text focuses on definitions and representative diagnostics at the minimal anchor; the material below provides additional evidence across refined balanced pairs $(n,m)$ and across bounded-complexity variant families.

\paragraph{Audit note (bounded scans, not free-form fitting).}
All tables in this appendix are produced by deterministic \emph{bounded scans} over explicitly specified finite families (loop sizes, phase denominators, phase-map families, and global relabelings), with deterministic tie-break rules.
The accompanying scripts write the resulting \LaTeX{} fragments for reproducibility; they do not introduce new degrees of freedom beyond the stated bounded families.

\subsection{Transport-rule sensitivity: padding/truncation and tie-breaks (audit)}
\label{subsec:holonomy_transport_rule_sensitivity}

The discrete edge transport of Section~\ref{subsec:deterministic_edge_connection} depends on three conventions: (i) how a Fold$_m$ fiber $P(w)$ is reduced to a fixed rank-$4$ representative list (truncation/selection and padding), and (ii) the deterministic tie-break used when multiple matchings attain the same minimum primary Hamming cost.
To rule out ``look-elsewhere'' effects from hidden degrees of freedom, we report a bounded counterfactual sensitivity envelope for gauge-invariant plaquette holonomy cycle statistics under two small families:
\begin{itemize}
\item \textbf{pad/select (H only):} vary only the rank-$4$ fiber representative (head/tail/spread selection and minimal padding), keeping the objective strictly Hamming;
\item \textbf{cost tie-break:} vary only a secondary tie-break among equal-Hamming minima, keeping padding/selection at the baseline.
\end{itemize}
\AuditTag The mainline construction fixes the tie-break as part of CAP (lexicographic minimality); the tie-break family is reported only as a transparent counterfactual sensitivity diagnostic, not as a tunable degree of freedom.
For each family we report (a) the fraction of 3/4-cycle plaquette holonomies $\mathrm{frac}_{3/4}$, (b) the maximum total-variation distance (TV) between the cycle-type distribution and the baseline distribution, and (c) the median minimum edge cost $q_{0.50}$ under the primary Hamming objective.

\begin{table}[t]
\centering
\scriptsize
\setlength{\tabcolsep}{3pt}
\renewcommand{\arraystretch}{1.15}
\resizebox{\textwidth}{!}{%
\begin{tabular}{llr l r r r r r r r r r l}
\toprule
$(n,m)$ & family & $|\mathcal{F}|$ & baseline rule &
$\mathrm{frac}_{3/4}$ & min & max & $\max|\Delta|$ & $\mathrm{TV}_{\max}$ &
$q_{0.50}$ & min & max & $\max|\Delta|$ & worst rule \\
\midrule
$(n,m)=(3,6)$ & \texttt{pad/select (H only)} & 6 & \texttt{head4/pad\_last/H} & 0.102 & 0.102 & 0.306 & 0.204 & 0.286 & 8 & 8 & 8 & 0 & \texttt{head4/pad\_first/H} \\
$(n,m)=(3,6)$ & \texttt{cost tie-break} & 3 & \texttt{head4/pad\_last/H} & 0.102 & 0.000 & 0.429 & 0.327 & 0.449 & 8 & 8 & 8 & 0 & \texttt{head4/pad\_last/H+wt} \\
$(n,m)=(4,8)$ & \texttt{pad/select (H only)} & 6 & \texttt{head4/pad\_last/H} & 0.231 & 0.120 & 0.249 & 0.111 & 0.151 & 9 & 9 & 10 & 1 & \texttt{spread4/pad\_first/H} \\
$(n,m)=(4,8)$ & \texttt{cost tie-break} & 3 & \texttt{head4/pad\_last/H} & 0.231 & 0.191 & 0.231 & 0.040 & 0.080 & 9 & 9 & 9 & 0 & \texttt{head4/pad\_last/H+wt} \\
\bottomrule%
\end{tabular}
}%
\caption{Sensitivity envelope of gauge-invariant plaquette holonomy cycle statistics under bounded counterfactual transport-rule families, for representative balanced anchors $(n,m)=(3,6)$ and $(4,8)$. Rows are generated by \texttt{scripts/exp\_holonomy\_transport\_rule\_sensitivity.py}.}
\label{tab:holonomy_transport_rule_sensitivity}
\end{table}

\subsection[Balanced-chain sweep]{Balanced-chain sweep across \texorpdfstring{$(n,m)=(1,2),(2,4),(3,6),(4,8),(5,10),(6,12),(7,14),(8,16)$}{(n,m)=(1,2),(2,4),(3,6),(4,8),(5,10),(6,12),(7,14),(8,16)}}
\label{subsec:holonomy_balanced_chain_sweep}

To probe how the finite connection/holonomy statistics behave under a minimal balanced refinement, we sweep the chain $m=2n$ for $n\in\{1,2,3,4,5,6,7,8\}$.
At each scale we build the same deterministic $S_4$-valued edge transport by truncating/padding each Fold$_m$ fiber to rank $4$ and selecting the minimum-cost bijection under Hamming distance on $m$-bit microstates.
We then summarize plaquette holonomy cycle types and the phase-lifted CP-odd signal at the phase denominator $\mathrm{denom}=2^m$.

\begin{table}[t]
\centering
\scriptsize
\setlength{\tabcolsep}{4pt}
\renewcommand{\arraystretch}{1.15}
\begin{tabular}{rrrrrrrrrrr}
\toprule
$n$ & $m$ & plaquettes &
1 & 2 & $2\times 2$ & 3 & 4 &
mean $|J|$ (3/4) & mean $J$ (3/4) & failures \\
\midrule
3 & 6 & 49 & 24 & 19 & 1 & 3 & 2 & 0.0470847 & +0.0207138 & 0 \\
4 & 8 & 225 & 126 & 44 & 3 & 42 & 10 & 0.0284115 & +0.0101108 & 0 \\
5 & 10 & 961 & 619 & 198 & 24 & 95 & 25 & 0.0325905 & +0.00800293 & 0 \\
\bottomrule%
\end{tabular}
\caption{Balanced-chain sweep for the holonomy and phase-lifted CP signal across $(n,m)\in\{(1,2),(2,4),(3,6),(4,8),(5,10),(6,12),(7,14),(8,16)\}$. Rows are generated by \texttt{scripts/exp\_holonomy\_balanced\_chain\_sweep.py}.}
\label{tab:holonomy_balanced_chain_sweep}
\end{table}

\paragraph{Balanced-chain permutation-robust mixing fits.}
Using the same phase-lifted holonomy extraction at each $(n,m)$ in the balanced chain, we can apply the same global $S_3\times S_3$ relabeling fit to PMNS- and CKM-style target sines.
Tables~\ref{tab:holonomy_balanced_chain_fit_pmns} and \ref{tab:holonomy_balanced_chain_fit_ckm} report the resulting best fits at each scale.

\begin{table}[t]
\centering
\scriptsize
\setlength{\tabcolsep}{4pt}
\renewcommand{\arraystretch}{1.15}
\begin{tabular}{rrrrrrrrrr}
\toprule
$n$ & $m$ & plaquettes & best $(\sigma_r,\sigma_c)$ & $s_{12}$ & $s_{23}$ & $s_{13}$ & $E_\infty$ & $E_1$ & mean $|J|$ \\
\midrule
3 & 6 & 49 & \texttt{(2, 0, 1)}/\texttt{(1, 2, 0)} & 0.7387 & 0.8513 & 0.1546 & 0.288 & 0.476 & 0.00480456 \\
4 & 8 & 225 & \texttt{(0, 2, 1)}/\texttt{(1, 0, 2)} & 0.8434 & 0.8706 & 0.1895 & 0.420 & 0.835 & 0.0068287 \\
5 & 10 & 961 & \texttt{(1, 2, 0)}/\texttt{(0, 1, 2)} & 0.8658 & 0.9017 & 0.1881 & 0.446 & 0.889 & 0.00420356 \\
\bottomrule%
\end{tabular}
\caption{Balanced-chain permutation-robust fit to PMNS target sines (finite diagnostic). Rows are generated by \texttt{scripts/exp\_holonomy\_balanced\_chain\_perm\_fit.py}.}
\label{tab:holonomy_balanced_chain_fit_pmns}
\end{table}

\begin{table}[t]
\centering
\scriptsize
\setlength{\tabcolsep}{4pt}
\renewcommand{\arraystretch}{1.15}
\begin{tabular}{rrrrrrrrrr}
\toprule
$n$ & $m$ & plaquettes & best $(\sigma_r,\sigma_c)$ & $s_{12}$ & $s_{23}$ & $s_{13}$ & $E_\infty$ & $E_1$ & mean $|J|$ \\
\midrule
3 & 6 & 49 & \texttt{(2, 1, 0)}/\texttt{(2, 1, 0)} & 0.4036 & 0.2930 & 0.1546 & 3.670 & 6.195 & 0.00480456 \\
4 & 8 & 225 & \texttt{(2, 1, 0)}/\texttt{(2, 1, 0)} & 0.2964 & 0.3150 & 0.1612 & 3.712 & 6.000 & 0.0068287 \\
5 & 10 & 961 & \texttt{(2, 1, 0)}/\texttt{(2, 1, 0)} & 0.2309 & 0.2424 & 0.1774 & 3.807 & 5.584 & 0.00420356 \\
\bottomrule%
\end{tabular}
\caption{Balanced-chain permutation-robust fit to CKM target sines (finite diagnostic). Rows are generated by \texttt{scripts/exp\_holonomy\_balanced\_chain\_perm\_fit.py}.}
\label{tab:holonomy_balanced_chain_fit_ckm}
\end{table}

\paragraph{Loop-scale sweep (k×k square holonomies).}
The unit-plaquette holonomy is the smallest closed loop.
As an additional finite-resolution diagnostic, we can compute phase-lifted holonomies around $k\times k$ square loops (for small $k$) and repeat the same 3/4-cycle restriction and $S_3\times S_3$ permutation-robust mixing fits.
Tables~\ref{tab:holonomy_loop_scale_cycle}--\ref{tab:holonomy_loop_scale_fit_ckm} summarize the resulting cycle-type counts and PMNS/CKM fit objectives for $k\in\{1,2,3\}$ on the $n=3$ grid.

\begin{table}[t]
\centering
\scriptsize
\setlength{\tabcolsep}{6pt}
\renewcommand{\arraystretch}{1.15}
\begin{tabular}{rrrrrrrr}
\toprule
$k$ & loops & 1 & 2 & $2\times 2$ & 3 & 4 & other \\
\midrule
1 & 49 & 24 & 19 & 1 & 3 & 2 & 0 \\
2 & 36 & 10 & 17 & 1 & 3 & 5 & 0 \\
3 & 25 & 7 & 10 & 1 & 6 & 1 & 0 \\
4 & 16 & 2 & 9 & 2 & 2 & 1 & 0 \\
5 & 9 & 1 & 3 & 0 & 5 & 0 & 0 \\
6 & 4 & 0 & 2 & 0 & 0 & 2 & 0 \\
7 & 1 & 0 & 0 & 0 & 0 & 1 & 0 \\
\bottomrule%
\end{tabular}
\caption{Cycle-type counts of $S_4$ holonomy permutations for $k\times k$ square loops on the $n=3$ grid. Rows are generated by \texttt{scripts/exp\_holonomy\_loop\_scale\_sweep.py}.}
\label{tab:holonomy_loop_scale_cycle}
\end{table}

\begin{table}[t]
\centering
\scriptsize
\setlength{\tabcolsep}{8pt}
\renewcommand{\arraystretch}{1.15}
\begin{tabular}{rrrrr}
\toprule
$k$ & 3/4 loops & mean angle [deg] & min [deg] & max [deg] \\
\midrule
1 & 5 & 108.000 & 90.000 & 120.000 \\
2 & 8 & 101.250 & 90.000 & 120.000 \\
3 & 7 & 115.714 & 90.000 & 120.000 \\
4 & 3 & 110.000 & 90.000 & 120.000 \\
5 & 5 & 120.000 & 120.000 & 120.000 \\
6 & 2 & 90.000 & 90.000 & 90.000 \\
7 & 1 & 90.000 & 90.000 & 90.000 \\
\bottomrule%
\end{tabular}
\caption{Loop-scale sweep of rotation angles in the sign-twisted standard $SO(3)\subset SU(3)$ representation bridge, restricted to 3/4-cycle holonomies. Rows are generated by \texttt{scripts/exp\_holonomy\_loop\_scale\_su3\_angle\_sweep.py}.}
\label{tab:holonomy_loop_scale_su3_angle}
\end{table}

\begin{table}[t]
\centering
\scriptsize
\setlength{\tabcolsep}{10pt}
\renewcommand{\arraystretch}{1.15}
\begin{tabular}{rrrrrr}
\toprule
$k$ & 3/4 loops & mean $W$ & min $W$ & max $W$ & mean $1-W$ \\
\midrule
1 & 5 & 0.318352 & 0.168535 & 0.428145 & 0.681648 \\
2 & 8 & 0.182942 & -0.00151624 & 0.428145 & 0.817058 \\
3 & 7 & 0.186643 & 0 & 0.377409 & 0.813357 \\
4 & 3 & 0.117196 & -0.0883019 & 0.441406 & 0.882804 \\
5 & 5 & 0.399583 & 0 & 0.642398 & 0.600417 \\
6 & 2 & 0.102449 & -0.0017954 & 0.206693 & 0.897551 \\
7 & 1 & -0.00151624 & -0.00151624 & -0.00151624 & 1.00152 \\
\bottomrule%
\end{tabular}
\caption{Loop-scale Wilson-loop style diagnostics $W=\Re(\mathrm{tr}(Q))/3$ from phase-lifted effective holonomies, restricted to 3/4-cycle loops \cite{Wilson1974,Rothe2012}. Rows are generated by \texttt{scripts/exp\_holonomy\_wilson\_loop\_sweep.py}.}
\label{tab:holonomy_wilson_loop}
\end{table}

\paragraph{Single-loop best fits (bounded scan).}
Finally, instead of averaging over a loop family, one can select a single loop together with a bounded phase denominator and a global $S_3\times S_3$ relabeling to best fit a target triple of sines.
Table~\ref{tab:holonomy_single_loop_bestfit} reports the best single-loop fits for PMNS- and CKM-style targets over a bounded search space.

\begin{table}[t]
\centering
\scriptsize
\setlength{\tabcolsep}{4pt}
\renewcommand{\arraystretch}{1.15}
\resizebox{\textwidth}{!}{%
\begin{tabular}{llrrrrlrrrrrrrr}
\toprule
target & map & denom & $k$ & $(x,y)$ & cycle & best $(\sigma_r,\sigma_c)$ & $s_{12}$ & $s_{23}$ & $s_{13}$ & $|J|$ & $E_\infty$ & $E_1$ & $E_\infty^{(2)}$ & $\Delta$ \\
\midrule
\texttt{PMNS} & \texttt{gray} & 64 & 4 & (1,0) & \texttt{2} & \texttt{(0, 1, 2)}/\texttt{(2, 0, 1)} & 0.6159 & 0.8334 & 0.164866 & 1.12757e-17 & 0.121 & 0.337 & 0.121 & 0.000 \\
\texttt{CKM} & \texttt{gray} & 512 & 3 & (3,0) & \texttt{2} & \texttt{(1, 0, 2)}/\texttt{(1, 0, 2)} & 0.1137 & 0.0804 & 0.00457281 & 1.79584e-19 & 0.679 & 1.473 & 0.706 & 0.027 \\
\bottomrule%
\end{tabular}
}%
\caption{Best single-loop fits to PMNS/CKM target sines under a bounded scan over $k\times k$ loops ($k\le 7$), phase denominators $\mathrm{denom}=2^p$ ($6\le p\le 18$), and global relabelings in $S_3\times S_3$. Rows are generated by \texttt{scripts/exp\_holonomy\_single\_loop\_bestfit.py}.}
\label{tab:holonomy_single_loop_bestfit}
\end{table}

\paragraph{Two-loop chains (bounded composition).}
One can also form a bounded family of effective holonomies by composing two selected loops (allowing inverses) before extracting the sines.
Table~\ref{tab:holonomy_two_loop_chain_bestfit} reports the best two-loop chain fits over a finite search family built from 3/4-cycle square loops.

\begin{table}[t]
\centering
\scriptsize
\setlength{\tabcolsep}{3pt}
\renewcommand{\arraystretch}{1.15}
\resizebox{\textwidth}{!}{%
\begin{tabular}{lrlrlllrrrrrrrr}
\toprule
target & $|\Theta|$ & map & denom & loop 1 & loop 2 & best $(\sigma_r,\sigma_c)$ & $s_{12}$ & $s_{23}$ & $s_{13}$ & $|J|$ & $E_\infty$ & $E_1$ & $E_\infty^{(2)}$ & $\Delta$ \\
\midrule
\texttt{PMNS} & 7195968 & \texttt{gray} & 256 & \texttt{(7,0,0,\texttt{4},-)} & \texttt{(2,5,4,\texttt{4},+)} & \texttt{(2,1,0)}/\texttt{(2,1,0)} & 0.5512 & 0.7422 & 0.149208 & 0.0262604 & 0.011 & 0.021 & 0.011 & 0.000 \\
\texttt{CKM} & 7195968 & \texttt{id} & 512 & \texttt{(6,0,1,\texttt{4},-)} & \texttt{(4,2,3,\texttt{3},+)} & \texttt{(1,2,0)}/\texttt{(2,1,0)} & 0.3490 & 0.0353 & 0.00287544 & 3.21877e-05 & 0.442 & 0.934 & 0.442 & 0.000 \\
\bottomrule%
\end{tabular}
}%
\caption{Best two-loop chain fits to PMNS/CKM target sines under a bounded scan over a finite family: (i) choose two square loops whose underlying $S_4$ holonomy is a 3/4-cycle, allowing inverses; (ii) choose a phase map (\texttt{id/gray/bitrev/not}) and denominator $\mathrm{denom}=2^p$ ($6\le p\le 18$); (iii) choose a global relabeling in $S_3\times S_3$. Rows are generated by \texttt{scripts/exp\_holonomy\_two\_loop\_chain\_bestfit.py}.}
\label{tab:holonomy_two_loop_chain_bestfit}
\end{table}

\paragraph{Two-loop chains with mixed cycle types.}
If one enlarges the admissible loop pool to include additional nontrivial holonomy cycle types (e.g.\ 2-cycles and $2\times 2$ cycles), one can obtain a substantially improved CKM-style fit within a still finite, auditable search box.
Table~\ref{tab:holonomy_two_loop_chain_mixed_cycles_bestfit} reports the best two-loop chain fits under a restricted phase family and a mixed-cycle loop pool.

\begin{table}[t]
\centering
\scriptsize
\setlength{\tabcolsep}{3pt}
\renewcommand{\arraystretch}{1.15}
\resizebox{\textwidth}{!}{%
\begin{tabular}{lrlrlllrrrrrrrr}
\toprule
target & $|\Theta|$ & map & denom & loop 1 & loop 2 & best $(\sigma_r,\sigma_c)$ & $s_{12}$ & $s_{23}$ & $s_{13}$ & $|J|$ & $E_\infty$ & $E_1$ & $E_\infty^{(2)}$ & $\Delta$ \\
\midrule
\texttt{PMNS} & 7962624 & \texttt{gray} & 256 & \texttt{(7,0,0,\texttt{4},-)} & \texttt{(2,5,4,\texttt{4},+)} & \texttt{(2,1,0)}/\texttt{(2,1,0)} & 0.5512 & 0.7422 & 0.149208 & 0.0262604 & 0.011 & 0.021 & 0.011 & 0.000 \\
\texttt{CKM} & 7962624 & \texttt{gray} & 512 & \texttt{(1,3,4,\texttt{2},+)} & \texttt{(2,3,5,\texttt{3},+)} & \texttt{(2,1,0)}/\texttt{(1,2,0)} & 0.2776 & 0.0374 & 0.00445891 & 3.40912e-05 & 0.213 & 0.459 & 0.213 & 0.000 \\
\bottomrule%
\end{tabular}
}%
\caption{Best two-loop chain fits to PMNS/CKM target sines under a bounded scan over mixed-cycle square loops (cycle types in \texttt{\{2,2x2,3,4\}}), with a restricted phase family (map \texttt{\{id,gray\}}, denom $\in\{256,512,1024\}$) and a global relabeling in $S_3\times S_3$. Rows are generated by \texttt{scripts/exp\_holonomy\_two\_loop\_chain\_mixed\_cycles\_bestfit.py}.}
\label{tab:holonomy_two_loop_chain_mixed_cycles_bestfit}
\end{table}

\begin{table}[t]
\centering
\scriptsize
\setlength{\tabcolsep}{4pt}
\renewcommand{\arraystretch}{1.15}
\begin{tabular}{rrrrrrrrrr}
\toprule
$k$ & loops & 3/4 loops & best $(\sigma_r,\sigma_c)$ & $s_{12}$ & $s_{23}$ & $s_{13}$ & $E_\infty$ & $E_1$ & mean $|J|$ \\
\midrule
1 & 49 & 5 & \texttt{(1, 2, 0)}/\texttt{(2, 0, 1)} & 0.5890 & 0.7457 & 0.3389 & 0.831 & 0.902 \\
2 & 36 & 8 & \texttt{(1, 2, 0)}/\texttt{(2, 0, 1)} & 0.5673 & 0.7137 & 0.3252 & 0.790 & 0.847 \\
3 & 25 & 7 & \texttt{(1, 0, 2)}/\texttt{(2, 0, 1)} & 0.4841 & 0.7469 & 0.3028 & 0.718 & 0.865 \\
\bottomrule%
\end{tabular}
\caption{Permutation-robust fit to PMNS target sines for $k\times k$ square holonomies (finite diagnostic), restricted to 3/4-cycle loops. Rows are generated by \texttt{scripts/exp\_holonomy\_loop\_scale\_sweep.py}.}
\label{tab:holonomy_loop_scale_fit_pmns}
\end{table}

\begin{table}[t]
\centering
\scriptsize
\setlength{\tabcolsep}{4pt}
\renewcommand{\arraystretch}{1.15}
\begin{tabular}{rrrrrrrrrr}
\toprule
$k$ & loops & 3/4 loops & best $(\sigma_r,\sigma_c)$ & $s_{12}$ & $s_{23}$ & $s_{13}$ & $E_\infty$ & $E_1$ & mean $|J|$ \\
\midrule
1 & 49 & 5 & \texttt{(1, 0, 2)}/\texttt{(2, 0, 1)} & 0.5890 & 0.6378 & 0.3389 & 4.455 & 8.136 & 0.0470847 \\
2 & 36 & 8 & \texttt{(1, 0, 2)}/\texttt{(2, 0, 1)} & 0.5673 & 0.6476 & 0.3252 & 4.413 & 8.072 & 0.0350064 \\
3 & 25 & 7 & \texttt{(1, 2, 0)}/\texttt{(2, 0, 1)} & 0.4841 & 0.6172 & 0.3028 & 4.342 & 7.794 & 0.0304645 \\
4 & 16 & 3 & \texttt{(2, 0, 1)}/\texttt{(1, 2, 0)} & 0.6599 & 0.6368 & 0.2988 & 4.329 & 8.122 & 0.0172157 \\
5 & 9 & 5 & \texttt{(2, 0, 1)}/\texttt{(2, 1, 0)} & 0.5003 & 0.6676 & 0.4070 & 4.638 & 8.201 & 0.0154748 \\
6 & 4 & 2 & \texttt{(0, 2, 1)}/\texttt{(1, 0, 2)} & 0.5569 & 0.3740 & 0.1825 & 3.836 & 6.927 & 0.0281961 \\
7 & 1 & 1 & \texttt{(2, 0, 1)}/\texttt{(1, 2, 0)} & 0.3556 & 0.5264 & 0.1623 & 3.718 & 6.703 & 0.0200472 \\
\bottomrule%
\end{tabular}
\caption{Permutation-robust fit to CKM target sines for $k\times k$ square holonomies (finite diagnostic), restricted to 3/4-cycle loops. Rows are generated by \texttt{scripts/exp\_holonomy\_loop\_scale\_sweep.py}.}
\label{tab:holonomy_loop_scale_fit_ckm}
\end{table}

\paragraph{Angle/CP trends under phase-denominator refinement.}
Finally, we can sweep the phase denominator $\mathrm{denom}=2^p$ and record the induced mean angles $(s_{12},s_{23},s_{13})$ and mean $|J|$ on the nontrivial holonomy subset (3/4 cycles), together with the log mismatch to $J_{\mathrm{geo}}$.
Table~\ref{tab:holonomy_phase_lift_angles_denom_sweep} reports this diagnostic sweep.

\begin{table}[t]
\centering
\scriptsize
\setlength{\tabcolsep}{6pt}
\renewcommand{\arraystretch}{1.15}
\begin{tabular}{rrrrrrrr}
\toprule
$\mathrm{denom}$ & $p$ & mean $s_{12}$ (3/4) & mean $s_{23}$ (3/4) & mean $s_{13}$ (3/4) & mean $|J|$ (3/4) & $\log(\mathrm{mean}/J_{\mathrm{geo}})$ & $|\cdot|$ \\
\midrule
64 & 6 & 0.6836 & 0.8197 & 0.4799 & 0.0470847 & +7.355 & 7.355 \\
128 & 7 & 0.6934 & 0.7867 & 0.4642 & 0.034465 & +7.043 & 7.043 \\
256 & 8 & 0.7156 & 0.8792 & 0.4045 & 0.0221376 & +6.601 & 6.601 \\
512 & 9 & 0.7186 & 0.8895 & 0.3655 & 0.010419 & +5.847 & 5.847 \\
1024 & 10 & 0.7194 & 0.8916 & 0.3463 & 0.00511243 & +5.135 & 5.135 \\
2048 & 11 & 0.7196 & 0.8921 & 0.3366 & 0.00254376 & +4.437 & 4.437 \\
4096 & 12 & 0.7196 & 0.8922 & 0.3316 & 0.00127032 & +3.743 & 3.743 \\
8192 & 13 & 0.7196 & 0.8923 & 0.3291 & 0.000634963 & +3.049 & 3.049 \\
16384 & 14 & 0.7196 & 0.8923 & 0.3279 & 0.000317457 & +2.356 & 2.356 \\
32768 & 15 & 0.7196 & 0.8923 & 0.3272 & 0.000158725 & +1.663 & 1.663 \\
65536 & 16 & 0.7196 & 0.8923 & 0.3269 & 7.93623e-05 & +0.970 & 0.970 \\
\textbf{131072} & 17 & 0.7196 & 0.8923 & 0.3268 & 3.96811e-05 & \textbf{+0.276} & \textbf{0.276} \\
262144 & 18 & 0.7196 & 0.8923 & 0.3267 & 1.98406e-05 & -0.417 & 0.417 \\
\bottomrule%
\end{tabular}
\caption{Phase-denominator sweep for mean extracted angles and mean $|J|$ on 3/4-cycle plaquettes at $n=3$ (finite diagnostic), together with the log mismatch to $J_{\mathrm{geo}}$. Rows are generated by \texttt{scripts/exp\_holonomy\_phase\_lift\_angles\_denom\_sweep.py}.}
\label{tab:holonomy_phase_lift_angles_denom_sweep}
\end{table}

\paragraph{A bounded-denominator PMNS fit (finite diagnostic).}
As an illustrative closure-style diagnostic, we can select $\mathrm{denom}=2^p$ to fit representative PMNS sines $(s_{12},s_{23},s_{13})$ using the same 3/4-cycle aggregated mean angles.
Table~\ref{tab:holonomy_pmns_denom_fit} reports the candidate sweep and the minimax/sum objectives.

\begin{table}[t]
\centering
\scriptsize
\setlength{\tabcolsep}{6pt}
\renewcommand{\arraystretch}{1.15}
\begin{tabular}{rrrrrrr}
\toprule
$\mathrm{denom}$ & $p$ & $s_{12}$ (3/4) & $s_{23}$ (3/4) & $s_{13}$ (3/4) & $E_\infty$ & $E_1$ \\
\midrule
64 & 6 & 0.6836 & 0.8197 & 0.4799 & 1.179 & 1.493 \\
128 & 7 & 0.6934 & 0.7867 & 0.4642 & 1.146 & 1.433 \\
256 & 8 & 0.7156 & 0.8792 & 0.4045 & 1.008 & 1.438 \\
512 & 9 & 0.7186 & 0.8895 & 0.3655 & 0.906 & 1.353 \\
1024 & 10 & 0.7194 & 0.8916 & 0.3463 & 0.853 & 1.302 \\
2048 & 11 & 0.7196 & 0.8921 & 0.3366 & 0.824 & 1.275 \\
4096 & 12 & 0.7196 & 0.8922 & 0.3316 & 0.809 & 1.260 \\
8192 & 13 & 0.7196 & 0.8923 & 0.3291 & 0.802 & 1.253 \\
16384 & 14 & 0.7196 & 0.8923 & 0.3279 & 0.798 & 1.249 \\
32768 & 15 & 0.7196 & 0.8923 & 0.3272 & 0.796 & 1.247 \\
65536 & 16 & 0.7196 & 0.8923 & 0.3269 & 0.795 & 1.246 \\
131072 & 17 & 0.7196 & 0.8923 & 0.3268 & 0.794 & 1.245 \\
\textbf{262144} & 18 & 0.7196 & 0.8923 & 0.3267 & \textbf{0.794} & \textbf{1.245} \\
\texttt{best/second} & $p=18/17$ & $-$ & $-$ & $-$ & 0.794/0.794 & $\Delta=0.000$ \\
\bottomrule%
\end{tabular}
\caption{Bounded-denominator fit to representative PMNS mixing sines using the phase-lifted holonomy angle extraction (finite diagnostic). Rows are generated by \texttt{scripts/exp\_holonomy\_phase\_lift\_pmns\_denom\_fit.py}.}
\label{tab:holonomy_pmns_denom_fit}
\end{table}

\paragraph{Permutation-robust fits (global relabeling search).}
The mapping from an effective $3\times 3$ unitary matrix to PDG angles depends on how row/column indices are identified with flavor/mass labels.
As a bounded-complexity diagnostic, we therefore allow a \emph{global} relabeling by a pair of permutations in $S_3\times S_3$ and select the best pair for each denominator by the same minimax objective on $(s_{12},s_{23},s_{13})$.
Tables~\ref{tab:holonomy_perm_fit_pmns} and \ref{tab:holonomy_perm_fit_ckm} report the resulting denominator sweeps for PMNS- and CKM-style target sines.

\begin{table}[t]
\centering
\scriptsize
\setlength{\tabcolsep}{4pt}
\renewcommand{\arraystretch}{1.15}
\begin{tabular}{rrrrrrrr}
\toprule
$\mathrm{denom}$ & $p$ & best $(\sigma_r,\sigma_c)$ & $s_{12}$ & $s_{23}$ & $s_{13}$ & $E_\infty$ & $E_1$ \\
\midrule
64 & 6 & \texttt{(1, 2, 0)}/\texttt{(2, 0, 1)} & 0.5890 & 0.7457 & 0.3389 & 0.831 & 0.902 \\
128 & 7 & \texttt{(1, 0, 2)}/\texttt{(2, 1, 0)} & 0.5263 & 0.7071 & 0.4583 & 1.133 & 1.227 \\
256 & 8 & \texttt{(1, 2, 0)}/\texttt{(2, 0, 1)} & 0.5279 & 0.6661 & 0.2715 & 0.609 & 0.760 \\
512 & 9 & \texttt{(1, 2, 0)}/\texttt{(2, 0, 1)} & 0.5273 & 0.6690 & 0.1979 & 0.293 & 0.441 \\
1024 & 10 & \texttt{(1, 2, 0)}/\texttt{(2, 0, 1)} & 0.5271 & 0.6693 & 0.1752 & 0.171 & 0.319 \\
2048 & 11 & \texttt{(1, 2, 0)}/\texttt{(2, 0, 1)} & 0.5270 & 0.6693 & 0.1688 & 0.134 & 0.282 \\
4096 & 12 & \texttt{(1, 2, 0)}/\texttt{(2, 0, 1)} & 0.5270 & 0.6693 & 0.1672 & 0.124 & 0.273 \\
8192 & 13 & \texttt{(1, 2, 0)}/\texttt{(2, 0, 1)} & 0.5270 & 0.6693 & 0.1668 & 0.122 & 0.270 \\
16384 & 14 & \texttt{(1, 2, 0)}/\texttt{(2, 0, 1)} & 0.5270 & 0.6693 & 0.1667 & 0.121 & 0.269 \\
32768 & 15 & \texttt{(1, 2, 0)}/\texttt{(2, 0, 1)} & 0.5270 & 0.6693 & 0.1667 & 0.121 & 0.269 \\
65536 & 16 & \texttt{(1, 2, 0)}/\texttt{(2, 0, 1)} & 0.5270 & 0.6693 & 0.1667 & 0.121 & 0.269 \\
131072 & 17 & \texttt{(1, 2, 0)}/\texttt{(2, 0, 1)} & 0.5270 & 0.6693 & 0.1667 & 0.121 & 0.269 \\
\textbf{262144} & 18 & \texttt{(1, 2, 0)}/\texttt{(2, 0, 1)} & 0.5270 & 0.6693 & 0.1667 & \textbf{0.121} & \textbf{0.269} \\
\texttt{best/second} & $-$ & $-$ & $-$ & $-$ & $-$ & 0.121/0.121 & $\Delta=0.000$ \\
\bottomrule%
\end{tabular}
\caption{Permutation-robust bounded-denominator fit to PMNS target sines using a global $S_3\times S_3$ relabeling. Rows are generated by \texttt{scripts/exp\_holonomy\_phase\_lift\_perm\_fit.py}.}
\label{tab:holonomy_perm_fit_pmns}
\end{table}

\begin{table}[t]
\centering
\scriptsize
\setlength{\tabcolsep}{4pt}
\renewcommand{\arraystretch}{1.15}
\begin{tabular}{rrrrrrrr}
\toprule
$\mathrm{denom}$ & $p$ & best $(\sigma_r,\sigma_c)$ & $s_{12}$ & $s_{23}$ & $s_{13}$ & $E_\infty$ & $E_1$ \\
\midrule
64 & 6 & \texttt{(1, 0, 2)}/\texttt{(2, 0, 1)} & 0.5890 & 0.6378 & 0.3389 & 4.455 & 8.136 \\
128 & 7 & \texttt{(1, 2, 0)}/\texttt{(2, 1, 0)} & 0.5263 & 0.6556 & 0.4583 & 4.756 & 8.352 \\
256 & 8 & \texttt{(1, 0, 2)}/\texttt{(2, 0, 1)} & 0.5279 & 0.6576 & 0.2715 & 4.233 & 7.835 \\
512 & 9 & \texttt{(1, 0, 2)}/\texttt{(2, 0, 1)} & 0.5273 & 0.6469 & 0.1979 & 3.916 & 7.501 \\
1024 & 10 & \texttt{(1, 0, 2)}/\texttt{(2, 0, 1)} & 0.5271 & 0.6448 & 0.1752 & 3.795 & 7.375 \\
2048 & 11 & \texttt{(1, 0, 2)}/\texttt{(2, 0, 1)} & 0.5270 & 0.6443 & 0.1688 & 3.758 & 7.338 \\
4096 & 12 & \texttt{(1, 0, 2)}/\texttt{(2, 0, 1)} & 0.5270 & 0.6441 & 0.1672 & 3.748 & 7.328 \\
8192 & 13 & \texttt{(1, 0, 2)}/\texttt{(2, 0, 1)} & 0.5270 & 0.6441 & 0.1668 & 3.746 & 7.325 \\
16384 & 14 & \texttt{(1, 0, 2)}/\texttt{(2, 0, 1)} & 0.5270 & 0.6441 & 0.1667 & 3.745 & 7.325 \\
32768 & 15 & \texttt{(1, 0, 2)}/\texttt{(2, 0, 1)} & 0.5270 & 0.6441 & 0.1667 & 3.745 & 7.324 \\
65536 & 16 & \texttt{(1, 0, 2)}/\texttt{(2, 0, 1)} & 0.5270 & 0.6441 & 0.1667 & 3.745 & 7.324 \\
131072 & 17 & \texttt{(1, 0, 2)}/\texttt{(2, 0, 1)} & 0.5270 & 0.6441 & 0.1667 & 3.745 & 7.324 \\
\textbf{262144} & 18 & \texttt{(1, 0, 2)}/\texttt{(2, 0, 1)} & 0.5270 & 0.6441 & 0.1667 & \textbf{3.745} & \textbf{7.324} \\
\texttt{best/second} & $-$ & $-$ & $-$ & $-$ & $-$ & 3.745/3.745 & $\Delta=0.000$ \\
\bottomrule%
\end{tabular}
\caption{Permutation-robust bounded-denominator fit to CKM target sines using a global $S_3\times S_3$ relabeling. Rows are generated by \texttt{scripts/exp\_holonomy\_phase\_lift\_perm\_fit.py}.}
\label{tab:holonomy_perm_fit_ckm}
\end{table}

\paragraph{Phase-map family sweep (low-complexity index transforms).}
The phase lift attaches a discrete phase to each microstate index.
To bound look-elsewhere freedom, we restrict to a small explicit family of low-complexity bit transforms $\tau$ (identity, Gray map, bit reversal, and complement) and rerun the bounded-denominator fits.
Tables~\ref{tab:holonomy_map_family_pmns} and \ref{tab:holonomy_map_family_ckm} report, for each map in a fixed small family, the best $(\mathrm{denom},\sigma_r,\sigma_c)$ under the same objective, together with the log mismatch of the resulting mean $|J|$ (3/4 cycles) to $J_{\mathrm{geo}}$.

\begin{table}[t]
\centering
\scriptsize
\setlength{\tabcolsep}{4pt}
\renewcommand{\arraystretch}{1.15}
\begin{tabular}{lrrlrrrrrr}
\toprule
map & $\mathrm{denom}$ & $p$ & best $(\sigma_r,\sigma_c)$ & $s_{12}$ & $s_{23}$ & $s_{13}$ & $E_\infty$ & $E_1$ & $\log(\mathrm{mean}|J|/J_{\mathrm{geo}})$ \\
\midrule
\texttt{id} & 262144 & 18 & \texttt{(1, 2, 0)}/\texttt{(2, 0, 1)} & 0.5270 & 0.6693 & 0.1667 & 0.121 & 0.269 & -0.417 \\
\texttt{gray} & 262144 & 18 & \texttt{(1, 2, 0)}/\texttt{(2, 0, 1)} & 0.5270 & 0.6693 & 0.1667 & 0.121 & 0.269 & +0.013 \\
\texttt{bitrev} & 262144 & 18 & \texttt{(1, 2, 0)}/\texttt{(2, 0, 1)} & 0.5270 & 0.6693 & 0.1667 & 0.121 & 0.269 & -0.245 \\
\texttt{not} & 262144 & 18 & \texttt{(1, 2, 0)}/\texttt{(2, 0, 1)} & 0.5270 & 0.6693 & 0.1667 & 0.121 & 0.269 & -0.417 \\
\bottomrule%
\end{tabular}
\caption{Phase-map family sweep for the PMNS target sines. Rows are generated by \texttt{scripts/exp\_holonomy\_phase\_lift\_map\_family\_sweep.py}.}
\label{tab:holonomy_map_family_pmns}
\end{table}

\begin{table}[t]
\centering
\scriptsize
\setlength{\tabcolsep}{4pt}
\renewcommand{\arraystretch}{1.15}
\begin{tabular}{lrrlrrrrrr}
\toprule
map & $\mathrm{denom}$ & $p$ & best $(\sigma_r,\sigma_c)$ & $s_{12}$ & $s_{23}$ & $s_{13}$ & $E_\infty$ & $E_1$ & $\log(\mathrm{mean}|J|/J_{\mathrm{geo}})$ \\
\midrule
\texttt{id} & 262144 & 18 & \texttt{(1, 0, 2)}/\texttt{(2, 0, 1)} & 0.5270 & 0.6441 & 0.1667 & 3.745 & 7.324 & -0.417 \\
\texttt{gray} & 262144 & 18 & \texttt{(1, 0, 2)}/\texttt{(2, 0, 1)} & 0.5270 & 0.6441 & 0.1667 & 3.745 & 7.324 & +0.013 \\
\texttt{bitrev} & 262144 & 18 & \texttt{(1, 0, 2)}/\texttt{(2, 0, 1)} & 0.5270 & 0.6441 & 0.1667 & 3.745 & 7.324 & -0.245 \\
\texttt{not} & 262144 & 18 & \texttt{(1, 0, 2)}/\texttt{(2, 0, 1)} & 0.5270 & 0.6441 & 0.1667 & 3.745 & 7.324 & -0.417 \\
\bottomrule%
\end{tabular}
\caption{Phase-map family sweep for the CKM target sines. Rows are generated by \texttt{scripts/exp\_holonomy\_phase\_lift\_map\_family\_sweep.py}.}
\label{tab:holonomy_map_family_ckm}
\end{table}

\subsection{A soft transport variant (temperature-like smoothing)}
\label{subsec:holonomy_soft_transport}

The strict minimum-cost matching rule produces a discrete $S_4$ transport on edges.
As an optional robustness diagnostic, one can form a \emph{soft} transport matrix from the full $4\times 4$ cost matrix between padded fibers, weighted by a temperature-like parameter $\beta$ via $\exp(-\beta\,\mathrm{cost})$, and then deterministically orthonormalize columns to obtain a unitary edge transport.
Sweeping $\beta$ interpolates between a highly mixed transport ($\beta\to 0$) and a sharp near-minimum transport (large $\beta$).
Tables~\ref{tab:holonomy_soft_transport_pmns} and \ref{tab:holonomy_soft_transport_ckm} report permutation-robust fits to PMNS- and CKM-style target sines under this soft transport, together with the mean $|J|$ on the resulting effective holonomies.

\begin{table}[t]
\centering
\scriptsize
\setlength{\tabcolsep}{4pt}
\renewcommand{\arraystretch}{1.15}
\begin{tabular}{rrrrrrrrrr}
\toprule
$\beta$ & plaquettes & best $(\sigma_r,\sigma_c)$ & $s_{12}$ & $s_{23}$ & $s_{13}$ & $E_\infty$ & $E_1$ & mean $|J|$ & failures \\
\midrule
0 & 25 & \texttt{(0, 1, 2)}/\texttt{(0, 2, 1)} & 0.5151 & 0.7661 & 0.4074 & 1.015 & 1.125 & 0.0246486 & 64 \\
0.25 & 25 & \texttt{(2, 0, 1)}/\texttt{(1, 2, 0)} & 0.6468 & 0.8308 & 0.4101 & 1.022 & 1.294 & 0.0351224 & 64 \\
\textbf{0.5} & 25 & \texttt{(2, 0, 1)}/\texttt{(2, 1, 0)} & 0.6458 & 0.8254 & 0.3966 & \textbf{0.988} & \textbf{1.253} & 0.0271959 & 64 \\
1 & 25 & \texttt{(2, 0, 1)}/\texttt{(2, 1, 0)} & 0.6283 & 0.7472 & 0.4718 & 1.162 & 1.299 & 0.0372322 & 64 \\
2 & 25 & \texttt{(0, 2, 1)}/\texttt{(0, 1, 2)} & 0.6554 & 0.6738 & 0.4933 & 1.206 & 1.466 & 0.0296337 & 64 \\
4 & 15 & \texttt{(0, 1, 2)}/\texttt{(1, 0, 2)} & 0.6508 & 0.7486 & 0.4525 & 1.120 & 1.295 & 0.0267211 & 93 \\
\bottomrule%
\end{tabular}
\caption{Soft-transport $\beta$ sweep with permutation-robust fit to PMNS target sines (robustness diagnostic). Rows are generated by \texttt{scripts/exp\_holonomy\_soft\_transport\_beta\_sweep.py}.}
\label{tab:holonomy_soft_transport_pmns}
\end{table}

\begin{table}[t]
\centering
\scriptsize
\setlength{\tabcolsep}{4pt}
\renewcommand{\arraystretch}{1.15}
\begin{tabular}{rrrrrrrrrr}
\toprule
$\beta$ & plaquettes & best $(\sigma_r,\sigma_c)$ & $s_{12}$ & $s_{23}$ & $s_{13}$ & $E_\infty$ & $E_1$ & mean $|J|$ & failures \\
\midrule
0 & 25 & \texttt{(0, 2, 1)}/\texttt{(0, 2, 1)} & 0.5151 & 0.5440 & 0.4074 & 4.639 & 8.026 & 0.0246486 & 64 \\
0.25 & 25 & \texttt{(2, 1, 0)}/\texttt{(1, 2, 0)} & 0.6468 & 0.5017 & 0.4101 & 4.645 & 8.180 & 0.0351224 & 64 \\
\textbf{0.5} & 25 & \texttt{(2, 1, 0)}/\texttt{(2, 1, 0)} & 0.6458 & 0.4799 & 0.3966 & \textbf{4.612} & \textbf{8.100} & 0.0271959 & 64 \\
1 & 25 & \texttt{(2, 1, 0)}/\texttt{(2, 1, 0)} & 0.6283 & 0.5826 & 0.4718 & 4.785 & 8.440 & 0.0372322 & 64 \\
2 & 25 & \texttt{(0, 1, 2)}/\texttt{(0, 1, 2)} & 0.6554 & 0.6457 & 0.4933 & 4.830 & 8.630 & 0.0296337 & 64 \\
4 & 15 & \texttt{(0, 2, 1)}/\texttt{(1, 0, 2)} & 0.6508 & 0.5703 & 0.4525 & 4.744 & 8.412 & 0.0267211 & 93 \\
\bottomrule%
\end{tabular}
\caption{Soft-transport $\beta$ sweep with permutation-robust fit to CKM target sines (robustness diagnostic). Rows are generated by \texttt{scripts/exp\_holonomy\_soft\_transport\_beta\_sweep.py}.}
\label{tab:holonomy_soft_transport_ckm}
\end{table}


